\documentclass[../../main.tex]{subfiles}
\begin{document}

\begin{center}
  {\fontsize{48pt}{58pt}\selectfont \textbf{1973}\normalsize\textbf{年度}}
\end{center}

\begin{table}[htb]
  \centering
  \caption{出題テーマと難易度}
  \begin{tabular}{|c|p{9cm}|c|}
    \hline
    \textbf{問題番号} & \textbf{テーマ} & \textbf{難易度} \\
    \hline
    第1問 & 三桁の数字を入れ替えて得る整数の集合 & \\
    \hline
    第2問 & 単位円上に根をもつ三次方程式 & \\
    \hline
    第3問 & 正三角形の位置ベクトルの関係 & \\
    \hline
    第4問 & 平方根を含む数列の極限 & \\
    \hline
    第5問 & 有限級数の値域の不等式 & \\
    \hline
    第6問 & くじを引く回数による期待報酬の比較 & \\
    \hline
  \end{tabular}
\end{table}

\begin{table}[htb]
  \centering
  \caption{解答}
  \renewcommand{\arraystretch}{1.4}
  \begin{tabular}{|c|c|p{9cm}|}
    \hline
    \textbf{問題番号} & \textbf{小問} & \textbf{答え} \\
    \hline
    第1問 & (1) & $A$は$19$個，$B$は$9$個 \\
    \cline{2-3}
    & (2) & $4455$ \\
    \hline
    第2問 & -- & $p=-q,r=-1,-1\le q\le3$，または$(p,q,r)=(-1,-1,1)$ \\
    \hline
    第3問 & -- & $\vec{c}=\dfrac{|\vec{b}|}{|\vec{a}|}\vec{a}+\dfrac{|\vec{a}|}{|\vec{b}|}\vec{b}$（証明問題） \\
    \hline
    第4問 & -- & $\dfrac{1}{4}$ \\
    \hline
    第5問 & -- & $-1\le x\le1$で$1\le f(x)<2$（証明問題） \\
    \hline
    第6問 & -- & $d=\dfrac{3-\sqrt{5}}{2}$とすると，$p<d$でA，$p=d$でA・B，$d<p<\dfrac{1}{2}$でB，$p=\dfrac{1}{2}$でB・C，$p>\dfrac{1}{2}$でCが最も有利 \\
    \hline
  \end{tabular}
\end{table}

\paragraph{試験形態}
6問．試験時間は要確認．

\end{document}
